A three-dimensional fluid can stretch its own rotation.
Follow a small group of the same fluid particles as they rotate around a common axis.
If the velocity field lengthens this material region along the axis, incompressibility forces the area of its cross-section to decrease.
When the rotation remains aligned with the extensional direction, the strain contributes positive amplification to the vorticity as the cross-section contracts.
The dangerous case is the one in which that geometry persists and the amplification survives the viscous response.
Then the next stretch would begin from a smaller cross-section and higher vorticity.
The problem is to decide whether the equation allows that repetition.

The rotation is measured by vorticity \(\omega\), and the deformation by the symmetric strain \(S\):
\[
  \omega:=\nabla\times u,
  \qquad
  S:=\frac12\bigl(\nabla u+(\nabla u)^{\mathsf T}\bigr).
\]
Taking the curl of \NavierStokes{} gives
\begin{equation}\label{eq:VorticityEquation}
  \partial_t\omega+(u\cdot\nabla)\omega
  =S\omega+\nu\Delta\omega.
\end{equation}
The left side follows the vorticity with the moving fluid.
The term \(S\omega\) is the aligned stretching contribution just described.
The term \(\nu\Delta\omega\) is viscous diffusion acting on the same spatial profile.
Both contributions act on the same vorticity during the same evolution.

At a smaller width, the next exact comparison is how \(S\omega\) and \(\nu\Delta\omega\) scale for one fixed normalized profile.
This isolates their homogeneities without assuming that the evolving vortex reproduces the profile or that scaling by itself decides the motion.
